% Complete Section 8.5, original PDF 111-112 / printed 101-102.
% block: parabolic-hyperbolic-heading page:111 type:section
\subsection{{Parabolic and hyperbolic problems}}\label{sec:parabolic-hyperbolic}

% block: hodge-heat-equation page:111 type:ordinary-paragraph
{Besides the Hodge Laplace equation $(d\delta+\delta d)u=f$, FEEC can be used to formulate and solve the related parabolic equation, which seeks a \emph{time-dependent} differential form $u:[0,T]\to\Lambda^k$ satisfying the \emph{Hodge heat equation}}
\begin{equation}\label{eq:hodge-heat-equation}
u_t+(\delta d+d\delta u)=f,
\end{equation}
{together with the \emph{initial condition} $u(0)=u_0$ and, for instance, the natural boundary conditions $\operatorname{tr}\star u=0$ and $\operatorname{tr}\star du=0$. As in the elliptic case, the weak formulation introduces the auxiliary variable $\sigma=\delta u$ and determines $(\sigma,u):[0,T]\to H\Lambda^{k-1}\times H\Lambda^k$ by}
\begin{equation}\label{eq:hodge-heat-weak-formulation}
\begin{aligned}
\langle\sigma,\tau\rangle-\langle d\tau,u\rangle&=0,&\tau&\in H\Lambda^{k-1},\ t\in[0,T],\\
\langle u_t,v\rangle+\langle d\sigma,v\rangle+\langle du,dv\rangle&=\langle f,v\rangle,&v&\in H\Lambda^k,\ t\in[0,T],
\end{aligned}
\end{equation}
{together with the initial condition. Note that the equations in \eqref{eq:hodge-heat-weak-formulation} must hold at each time $t$. In comparison to \eqref{eq:general-hodge-weak-formulation}, \eqref{eq:hodge-heat-weak-formulation} includes an additional term involving the time-derivative $u_t$, but, unlike the elliptic case, even if there are harmonic forms, they do not need to be explicitly accounted for in the weak formulation of the parabolic problem.}

% block: hodge-heat-analysis pages:111,112 type:ordinary-paragraph
{The Hodge heat equation is discussed in the FEEC setting in [\BookNumber{10}, Theorem 4.4]. There we show how the results we have obtained for the Hodge Laplacian can be used together with the Hille–Yosida–Phillips theory to establish well-posedness of the parabolic problem. The paper goes on to analyze the Galerkin method, both semidiscrete and, using a simple time-stepping scheme, fully discrete. Optimal order error estimates are obtained based on an appropriate elliptic projection, whose properties again follow from the elliptic case.}

% block: hodge-wave-equation page:112 type:ordinary-paragraph
{We can also consider the \emph{Hodge wave equation}, which is a second order hyperbolic equation determining a time-dependent $k$-form $u:[0,T]\to\Lambda^k$ from the differential equation}
\begin{equation}\label{eq:hodge-wave-equation}
u_{tt}+(\delta d+d\delta)u=f.
\end{equation}

% block: hodge-wave-first-order-system page:112 type:ordinary-paragraph
{Of course, we now require two initial conditions, $u(0)=u_0$ and $u_t(0)=v_0$. To solve this problem we introduce three auxiliary variables: $\sigma=\delta u$, $v=u_t$, and $\beta=du$, time-dependent differential forms of order $k-1$, $k$, and $k+1$, respectively. We may then derive the strong form equations}
\[
\begin{aligned}
\sigma_t&=\delta u_t=\delta v,\\
v_t&=u_{tt}=-(\delta d+d\delta)u+f=-\delta\beta-d\sigma+f,\\
\beta_t&=du_t=dv,
\end{aligned}
\]
{or, in matrix form,}
\[
\frac d{dt}\begin{pmatrix}\sigma\\v\\\beta\end{pmatrix}=
\begin{pmatrix}0&\delta&0\\-d&0&-\delta\\0&d&0\end{pmatrix}
\begin{pmatrix}\sigma\\v\\\beta\end{pmatrix}+\begin{pmatrix}0\\f\\0\end{pmatrix},
\]
{which reveals the structure of a symmetric hyperbolic system (the key point being that the matrix operator on the right-hand side is skew-symmetric). As initial data, we set $\sigma(0)=\delta u_0$, $v(0)=v_0$, and $\beta(0)=du_0$. Note that the original variable $u$ is not part of the system but can be obtained by integrating $u_t=v$, i.e.,}
\[
u(t)=u_0+\int_0^t v(s)\,ds.
\]

% block: hodge-wave-weak-formulation page:112 type:ordinary-paragraph
{Turning now to the mixed weak formulation, it seeks $(\sigma,v,\beta):[0,T]\to H\Lambda^{k-1}\times H\Lambda^k\times H\Lambda^{k+1}$ satisfying}
\begin{equation}\label{eq:hodge-wave-weak-formulation}
\begin{aligned}
\langle\sigma_t,\tau\rangle-\langle v,d\tau\rangle&=0,&\tau&\in H\Lambda^{k-1},\ t\in[0,T],\\
\langle v_t,w\rangle+\langle d\sigma,w\rangle+\langle\beta,dw\rangle&=\langle f,w\rangle,&w&\in H\Lambda^k,\ t\in[0,T],\\
\langle\beta_t,\gamma\rangle-\langle dv,\gamma\rangle&=0,&\gamma&\in H\Lambda^{k+1},\ t\in[0,T].
\end{aligned}
\end{equation}

% block: hodge-wave-wellposedness-energy page:112 type:ordinary-paragraph
{The well-posedness of this formulation and the convergence of its Galerkin solution are established in the Hilbert complex framework in [\BookNumber{63}, Chapter 4]. As for the parabolic problem, this is obtained using Hille–Yosida–Phillips theory, which is particularly simple when, as here, the time-independent operator is skew-symmetric. The same skew-symmetry leads immediately to energy conservation, as well. Namely, taking the test functions $(\tau,w,\gamma)$ equal to the solution $(\sigma,v,\beta)$ at time $t$ in the weak formulation, we see that}
\[
\|\sigma\|_{L^2}^2+\|v\|_{L^2}^2+\|\beta\|_{L^2}^2
\]
{remains constant in time.}
